\clearpage

\setcounter{chapter}{6}
\setcounter{section}{0}

% Add the appendices to table of contents
\addcontentsline{toc}{chapter}{\numberline{A}Appendices}

% Set up page style for appendices
\pagestyle{fancy}
\fancyhf{} % Clear all header and footer fields
\fancyhead[L]{\footnotesize\textit{Ars Post Faber: Digital fabrication democratization through embodied knowledge preservation}}
\fancyfoot[C]{\thepage} % Page number in footer center
\renewcommand{\headrulewidth}{0pt}
\renewcommand{\footrulewidth}{0pt}

% Custom appendices title to match chapter formatting
\noindent
{\Large\textbf{Appendices}}
\vspace{0.3cm}
\hrule
\vspace{0.8cm}
\label{ch:appendices}

% Set no paragraph indentation
\setlength{\parindent}{0pt}

\clearpage
\section{\textit{AI.RTISANSHIP} technical documentation}

The \textit{AI.RTISANSHIP} experiment represented an attempt to capture and digitize traditional pottery techniques through computer vision and machine learning. This technical documentation outlines the complete implementation pipeline, from data collection through real-time prediction deployment.

\vspace{0.5cm}

The code for this project can be found opened under an MIT license on the following GitHub repository: \href{https://github.com/jmuozan/ML_Sequence_Recognition}{github.com/jmuozan/ML\_Sequence\_Recognition}

\subsection{System architecture overview}

The system implemented a complete machine learning pipeline organized into seven main sections within a Jupyter notebook environment (\texttt{ML\_Sequence\_Recognition.ipynb}), containing 53 cells that progressively built from data collection to real-time prediction capabilities.

\vspace{0.5cm}

\begin{enumerate}
\item Data capture using MediaPipe holistic landmark detection
\item Manual labeling and quality assurance of movement sequences
\item Data preprocessing and feature extraction
\item Neural network training with bidirectional LSTM architecture
\item Model validation and performance evaluation
\item Real-time prediction system deployment
\end{enumerate}

\vspace{0.5cm}

\subsection{Data collection and mediaPipe integration}

\subsubsection{Library installation and setup}

The system began with installation of core computer vision and machine learning dependencies:

\begin{verbatim}
pip install mediapipe opencv-python pandas numpy matplotlib
\end{verbatim}

Essential library imports establishing the environment:

\begin{verbatim}
import mediapipe as mp
import cv2
import numpy as np
import time
import pandas as pd
import csv
import os
from matplotlib import pyplot as plt
import math
\end{verbatim}

\subsubsection{MediaPipe holistic model configuration}

The landmark detection system used MediaPipe's holistic model, configured for simultaneous pose and hand tracking:

\begin{verbatim}
mp_pose = mp.solutions.pose
mp_drawing = mp.solutions.drawing_utils

cap = cv2.VideoCapture(0)
with mp_pose.Pose(min_detection_confidence=0.5, 
                  min_tracking_confidence=0.5) as pose, \
    mp.solutions.hands.Hands(min_detection_confidence=0.5, 
                            min_tracking_confidence=0.5) as hands:
    
    while cap.isOpened():
        ret, frame = cap.read()
        if not ret:
            break
            
        # Convert BGR to RGB
        image = cv2.cvtColor(frame, cv2.COLOR_BGR2RGB)
        image.flags.writeable = False
        
        # Make detection
        pose_results = pose.process(image)
        hand_results = hands.process(image)
        
        # Convert back to BGR for display
        image.flags.writeable = True
        image = cv2.cvtColor(image, cv2.COLOR_RGB2BGR)
\end{verbatim}

Where the detection confidence thresholds of 0.5 looked to find certain balance between accuracy and speed, while the tracking confidence maintained certain consistency consistency across frames.

\subsubsection{Feature vector construction}

After, the system extracts 225-dimensional feature vectors from each frame, combining pose and hand landmark coordinates:

\begin{verbatim}
def extract_keypoints(results):
    # Extract pose landmarks (33 points x 3 coordinates = 99 features)
    pose = np.array([[res.x, res.y, res.z] for res in 
                     results.pose_landmarks.landmark]).flatten() \
           if results.pose_landmarks else np.zeros(33*3)
    
    # Extract left hand landmarks (21 points x 3 coordinates = 63 features)
    lh = np.array([[res.x, res.y, res.z] for res in 
                   results.left_hand_landmarks.landmark]).flatten() \
         if results.left_hand_landmarks else np.zeros(21*3)
    
    # Extract right hand landmarks (21 points x 3 coordinates = 63 features)
    rh = np.array([[res.x, res.y, res.z] for res in 
                   results.right_hand_landmarks.landmark]).flatten() \
         if results.right_hand_landmarks else np.zeros(21*3)
    
    return np.concatenate([pose, lh, rh])
\end{verbatim}

\subsubsection{Data collection workflow}

In order to process the data collection, this script implemented a semi-automated labeling through keyboard interaction:

\begin{verbatim}
# CSV header structure for 225 features + metadata
headers = ['class', 'accuracy', 'sequence'] 
headers.extend(['pose_' + coord + str(i) for i in range(33) 
                for coord in ('x', 'y', 'z')])
headers.extend([hand + '_' + coord + str(i) for hand in ('right_hand', 'left_hand') 
                for i in range(21) for coord in ('x', 'y', 'z')])

# Recording control system
key = cv2.waitKey(10) & 0xFF
if key == ord('r'):  # Start recording correct sequence
    recording = True
    current_label = 'correct'
elif key == ord('w'):  # Start recording incorrect sequence
    recording = True
    current_label = 'incorrect'
elif key == ord('s'):  # Stop recording
    recording = False
elif key == ord('q'):  # Quit application
    break
\end{verbatim}

\subsection{Data processing and feature extraction}

\subsubsection{Consolidation}

Multiple CSV files from different recording sessions and movements were combined into a unified training dataset using pandas:

\begin{verbatim}
# Combine multiple data collection sessions
df1 = pd.read_csv('coordinates_1.csv', header=None)
df2 = pd.read_csv('coordinates_2.csv', header=None) 
combined_df = pd.concat([df1, df2], axis=0, ignore_index=True)

# Verify feature dimensions
rows, columns = combined_df.shape
print("Dataset shape: " + str(rows) + " rows, " + str(columns) + " columns")
print("Expected features: " + str(21*3 + 21*3 + 33*3) + " = 225")
\end{verbatim}

\subsubsection{Labeling}

The labeling system combined movement type with correctness assessment (asigning "true" if correct and "flase" if incorrect):

\begin{verbatim}
# Transform binary accuracy to readable labels
data['accuracy'] = data['accuracy'].map({0: 'W', 1: 'R'})

# Create composite labels (e.g., "pottery_R", "pottery_W")
data['label'] = data.iloc[:, 0].astype(str) + '_' + data['accuracy']

# Extract unique class labels
unique_labels = data['label'].unique()
print("Classes identified: " + str(unique_labels))
\end{verbatim}

\subsubsection{Sequence organization}

This temporal labeled data was organized into sequence-based directory structures for the posterior LSTM training:

\begin{verbatim}
def create_sequence_directories(csv_file, base_folder):
    """
    Organize frame data into sequence-based directory structure
    """
    data = pd.read_csv(csv_file)
    
    for label in data['label'].unique():
        label_folder = os.path.join(base_folder, label)
        os.makedirs(label_folder, exist_ok=True)
        
        label_data = data[data['label'] == label]
        sequences = label_data['sequence'].unique()
        
        for seq_num in sequences:
            seq_folder = os.path.join(label_folder, str(seq_num))
            os.makedirs(seq_folder, exist_ok=True)
            
            seq_data = label_data[label_data['sequence'] == seq_num]
            for idx, row in seq_data.iterrows():
                features = row.iloc[3:].values  # Skip metadata columns
                np.save(os.path.join(seq_folder, str(idx) + ".npy"), features)
\end{verbatim}

\subsection{Neural network architecture}

\subsubsection{Model design}

The classification system implemented a three-layer bidirectional LSTM architecture optimized for temporal sequence recognition:

\begin{verbatim}
from tensorflow.keras.models import Sequential
from tensorflow.keras.layers import LSTM, Dense, Dropout, Bidirectional
from tensorflow.keras.regularizers import l2
from tensorflow.keras.optimizers import Adam

# Build bidirectional LSTM architecture
model = Sequential([
    # First bidirectional layer with return sequences
    Bidirectional(LSTM(64, return_sequences=True, activation='tanh'), 
                  input_shape=(sequence_length, 225)),
    Dropout(0.2),
    
    # Second bidirectional layer with L2 regularization
    Bidirectional(LSTM(128, return_sequences=True, activation='tanh', 
                       kernel_regularizer=l2(0.01))),
    Dropout(0.2),
    
    # Final bidirectional layer
    Bidirectional(LSTM(64, activation='tanh')),
    
    # Dense classification layers
    Dense(64, activation='relu'),
    Dense(32, activation='relu'),
    Dense(num_classes, activation='softmax')
])

# Compile model with Adam optimizer
optimizer = Adam(learning_rate=0.001)
model.compile(optimizer=optimizer, 
              loss='categorical_crossentropy', 
              metrics=['accuracy'])
\end{verbatim}

\subsubsection{Training configuration}

Training callbacks were implemented for learning control:

\begin{verbatim}
from tensorflow.keras.callbacks import EarlyStopping, ModelCheckpoint, ReduceLROnPlateau

# Configure training callbacks
early_stopping = EarlyStopping(monitor='val_accuracy', 
                               patience=20, 
                               restore_best_weights=True)

checkpoint = ModelCheckpoint('best_model.keras', 
                            save_best_only=True, 
                            monitor='val_accuracy', 
                            mode='max')

reduce_lr = ReduceLROnPlateau(monitor='val_loss', 
                             factor=0.2, 
                             patience=5, 
                             min_lr=0.0001)

# Execute training with monitoring
history = model.fit(X_train, y_train, 
                   validation_split=0.2, 
                   epochs=2000, 
                   batch_size=32,
                   callbacks=[early_stopping, checkpoint, reduce_lr])
\end{verbatim}

\subsection{Real-time prediction}

\subsubsection{Live processing pipeline}

In order to test the proper functioning of 'actions.h5' a real-time system integrating the MediaPipe detection with trained model inference was implemented:

\begin{verbatim}
def mediapipe_detection(image, model):
    # Process frame through MediaPipe pipeline
    image = cv2.cvtColor(image, cv2.COLOR_BGR2RGB)
    image.flags.writeable = False
    results = model.process(image)
    image.flags.writeable = True
    image = cv2.cvtColor(image, cv2.COLOR_RGB2BGR)
    return image, results

def draw_styled_landmarks(image, results, incorrect_movement):
    # Render landmarks with correctness-based color coding
    color = (0, 0, 255) if incorrect_movement else (0, 255, 0)  # Red/Green
    
    # Draw pose landmarks
    if results.pose_landmarks:
        mp_drawing.draw_landmarks(
            image, results.pose_landmarks, mp_pose.POSE_CONNECTIONS,
            mp_drawing.DrawingSpec(color=color, thickness=2),
            mp_drawing.DrawingSpec(color=color, thickness=2))
    
    # Draw hand landmarks
    if results.left_hand_landmarks:
        mp_drawing.draw_landmarks(
            image, results.left_hand_landmarks, mp_hands.HAND_CONNECTIONS,
            mp_drawing.DrawingSpec(color=color, thickness=2),
            mp_drawing.DrawingSpec(color=color, thickness=2))
\end{verbatim}

\subsubsection{Rolling buffer implementation}

The system designed looked to maintain certain temporal context to be able to make the predictions through a frame buffer, saving a certain amount of frames and processign constantly:

\begin{verbatim}
# Initialize system components
sequence = []
cap = cv2.VideoCapture(0)
model = load_model('pottery_model.h5')

with mp_holistic.Holistic(min_detection_confidence=0.5, 
                          min_tracking_confidence=0.5) as holistic:
    while cap.isOpened():
        ret, frame = cap.read()
        frame = cv2.flip(frame, 1)  # Mirror image
        
        # Process frame
        image, results = mediapipe_detection(frame, holistic)
        
        # Extract and buffer keypoints
        if results.pose_landmarks or results.left_hand_landmarks or results.right_hand_landmarks:
            keypoints = extract_keypoints(results)
            sequence.append(keypoints)
            sequence = sequence[-30:]  # Maintain 30-frame rolling window
            
            # Predict when buffer is full
            if len(sequence) == 30:
                res = model.predict(np.expand_dims(sequence, axis=0))[0]
                action = actions[np.argmax(res)]
                incorrect_movement = action.endswith("_W")
                
                # Visual feedback
                draw_styled_landmarks(image, results, incorrect_movement)
                color = (0, 0, 255) if incorrect_movement else (0, 255, 0)
                cv2.putText(image, 'Action: ' + str(action), (15, 50), 
                           cv2.FONT_HERSHEY_SIMPLEX, 1, color, 2)
        
        cv2.imshow('AI.RTISANSHIP', image)
        if cv2.waitKey(10) & 0xFF == ord('q'):
            break
\end{verbatim}

\subsection{Performance evaluation}

\subsubsection{Metrics}

Metrics were extracted to see if the trained model achieved reasonable accuracy in distinguishing between predefined movement categories:

\begin{verbatim}
from sklearn.metrics import multilabel_confusion_matrix, accuracy_score

# Generate predictions on test set
yhat = model.predict(X_test)
ytrue = np.argmax(y_test, axis=1).tolist()  
yhat = np.argmax(yhat, axis=1).tolist()

# Calculate accuracy metrics
accuracy = accuracy_score(ytrue, yhat)
print("Test Accuracy: " + str(round(accuracy, 3)))

# Generate confusion matrix
confusion_matrices = multilabel_confusion_matrix(ytrue, yhat)
for i, matrix in enumerate(confusion_matrices):
    print("Class " + str(actions[i]) + " Confusion Matrix:")
    print(matrix)
\end{verbatim}

\subsection{Technical limitations}

The implementation revealed several critical limitations that informed the broader research conclusions:

\subsubsection{Sensory constraints}
Computer vision systems remained blind to essential tactile feedback channels that pottery practice depends upon, including clay resistance, texture variations, and vibrational cues.

\subsubsection{Individual variation}
Each participating artisan attempted different approaches to identical tasks, reflecting personal histories and preferences that "resist" standardization into training categories.

\subsubsection{Adaptive vs. pattern-based practice}
The machine learning system's requirement for discrete classification categories proved incompatible with pottery's emphasis on adaptation to emergent material conditions.

\clearpage
\section{\textit{CR3ATED} Technical documentation}


\clearpage
\section{\textit{AI Tools} Technical documentation}


\clearpage
\section{\textit{Ars Post Faber} Technical documentation}
